% from ega1-6-fr.tex
\oldpage[I]{141}
\emph{Module (resp. tout $\mathscr{O}_X$-Module quotient) quasi-cohérent}
d'un $\mathscr{O}_X$-Module \emph{cohérent} $\mathscr{F}$ est alors
\emph{cohérent}, car la question est de nouveau locale, et il suffit
d'appliquer (\hyperref[I.1.5.1-fr]{1.5.1}),
(\hyperref[I.1.4.1-fr]{1.4.1}) et
(\hyperref[I.1.3.10-fr]{1.3.10}), joints au fait qu'un sous-module (resp.
module quotient) d'un module de type fini sur un anneau noethérien est de
type fini. Plus particulièrement, \emph{tout faisceau d'idéaux
quasi-cohérent} de $\mathscr{O}_X$ est \emph{cohérent}.

Si un préschéma $X$ est réunion finie (resp. réunion) d'ouverts $W_\lambda$
tels que les préschémas induits sur les $W_\lambda$ soient noethériens (resp.
localement noethériens), il est clair que $X$ est noethérien (resp.
localement noethérien).

\begin{proposition}[6.1.2]
\label{I.6.1.2-fr}
Pour qu'un préschéma $X$ soit noethérien, il faut et il suffit qu'il soit
localement noethérien et que son espace sous-jacent soit quasi-compact.
L'espace sous-jacent de $X$ est alors noethérien.
\end{proposition}

\begin{proof}
La première assertion découle aussitôt des définitions et de
(\hyperref[I.1.1.10-fr]{1.1.10 (ii)}). La seconde résulte de
(\hyperref[I.1.1.6-fr]{1.1.6}) et du fait que tout espace réunion finie de
sous-espaces noethériens est noethérien
(\hyperref[0.2.2.3-fr]{0, 2.2.3}).
\end{proof}

\begin{proposition}[6.1.3]
\label{I.6.1.3-fr}
Soit $X$ un schéma affine d'anneau $A$. Les conditions suivantes sont
équivalentes~: a) $X$ est noethérien~; b) $X$ est localement noethérien~;
c) $A$ est noethérien.
\end{proposition}

\begin{proof}
L'équivalence de a) et b) résulte de (\hyperref[I.6.1.2-fr]{6.1.2}) et de ce
que tout schéma affine a un espace sous-jacent quasi-compact
(\hyperref[I.1.1.10-fr]{1.1.10})~; il est clair en outre que c) entraîne a).
Pour voir que a) entraîne c), remarquons qu'il y a un recouvrement fini
$(V_i)$ de $X$ par des ouverts affines tels que l'anneau $A_i$ du préschéma
induit sur $V_i$ soit noethérien. Soit alors $(\mathfrak{a}_n)$ une suite
croissante d'idéaux de $A$~; il lui correspond canoniquement de façon
biunivoque (\hyperref[I.1.3.7-fr]{1.3.7}) une suite croissante
$(\widetilde{\mathfrak{a}}_n)$ de faisceaux d'idéaux dans
$\widetilde{A}\equiv\mathscr{O}_X$~; pour voir que la suite
$(\mathfrak{a}_n)$ est stationnaire, il suffit de prouver que la suite
$(\widetilde{\mathfrak{a}}_n)$ l'est. Or, la restriction
$\widetilde{\mathfrak{a}}_n|V_i$ est un faisceau quasi-cohérent d'idéaux dans
$\mathscr{O}_X|V_i$, étant l'image réciproque de
$\widetilde{\mathfrak{a}}_n$ par l'injection canonique $V_i\to X$
(\hyperref[0.5.1.4-fr]{0, 5.1.4})~; $\widetilde{\mathfrak{a}}_n|V_i$ est
donc de la forme $\widetilde{\mathfrak{a}}_{ni}$, où $\mathfrak{a}_{ni}$ est
un idéal de $A_i$ (\hyperref[I.1.3.7-fr]{1.3.7}). Comme $A_i$ est noethérien,
la suite $(\mathfrak{a}_{ni})$ est stationnaire pour tout $i$, d'où la
proposition.
\end{proof}

On notera que le raisonnement précédent prouve aussi que \emph{si $X$ est un
préschéma noethérien, toute suite croissante de faisceaux cohérents d'idéaux
de $\mathscr{O}_X$ est stationnaire}.

\begin{proposition}[6.1.4]
\label{I.6.1.4-fr}
Tout sous-préschéma d'un préschéma noethérien (resp. localement noethérien)
est noethérien (resp. localement noethérien).
\end{proposition}

\begin{proof}
Il suffit de faire la démonstration pour un préschéma noethérien $X$~; en
outre, on est aussitôt ramené par la définition
(\hyperref[I.6.1.1-fr]{6.1.1}) au cas où $X$ est un schéma affine. Comme tout
sous-préschéma de $X$ est un sous-préschéma fermé d'un préschéma induit sur
un ouvert (\hyperref[I.4.1.3-fr]{4.1.3}), on peut se borner à considérer le
cas d'un sous-préschéma $Y$ fermé ou induit sur un ouvert de $X$. Le cas où
$Y$ est fermé est immédiat, car si $A$ est l'anneau de $X$, on sait que $Y$
est un schéma affine d'anneau $A/\mathfrak{J}$, où $\mathfrak{J}$ est un
idéal de $A$ (\hyperref[I.4.2.3-fr]{4.2.3})~; comme $A$ est noethérien
(\hyperref[I.6.1.3-fr]{6.1.3}), il en est de même de $A/\mathfrak{J}$.

Supposons maintenant $Y$ ouvert dans $X$~; l'espace sous-jacent $Y$ est
noethérien (\hyperref[I.6.1.2-fr]{6.1.2}), donc quasi-compact, et par suite
réunion finie d'ouverts $D(f_i)$ ($f_i\in A$)~;
